% Clase 8 — el metodo de las imagenes, caso plano. La carga imagen es FICTICIA y
% la solucion vale SOLO del lado de la carga real: del otro lado del espejo es
% otra historia.
% Las dos cotas `d` van bien abajo y separadas: a la altura del eje se pisan entre
% si y con los rotulos de las cargas.
\begin{tikzpicture}[scale=1]

  % region donde NO vale la solucion
  \fill[figgray!12] (-3.3,-2.6) rectangle (0,2.2);
  \draw[figline, line width=1.4pt] (0,-2.6) -- (0,2.2);
  \node[figlbl, figblue] at (0,2.5) {plano conductor};
  \node[figlblaux, align=center] at (-2.55,-1.35) {aqu\'i la soluci\'on\\ \emph{no} vale};

  \draw[figaxis] (-3.3,0) -- (3.5,0);
  \node[figlblaux] at (3.7,0) {$x$};

  % carga real
  \node[figpos] at (1.7,0) {};
  \node[figlbl, figred] at (1.7,0.42) {$+q$};

  % carga imagen
  \node[figneg] at (-1.7,0) {};
  \node[figlbl, figblue] at (-1.7,0.42) {$q'=-q$};
  \node[figlblaux, align=center] at (-1.7,1.28) {imagen\\ (ficticia)};

  % carga inducida real sobre el plano
  \foreach \yy in {-1.2,-0.85,0.85,1.2}{\node[figneg] at (0,\yy) {};}
  \node[figlblaux, figblue, fill=white, inner sep=1.2pt] at (0.85,1.55) {$\sigma<0$};

  % cotas, bien abajo
  \draw[figaux] (-1.7,-0.22) -- (-1.7,-2.15);
  \draw[figaux] (1.7,-0.22) -- (1.7,-2.15);
  \draw[figaux] (0,-1.45) -- (0,-2.15);
  \draw[figvecaux] (0,-1.95) -- (1.7,-1.95);
  \draw[figvecaux] (1.7,-1.95) -- (0,-1.95);
  \node[figlbl, fill=white, inner sep=1.5pt] at (0.85,-1.95) {$d$};
  \draw[figvecaux] (0,-1.95) -- (-1.7,-1.95);
  \draw[figvecaux] (-1.7,-1.95) -- (0,-1.95);
  \node[figlbl, fill=figgray!12, inner sep=1.5pt] at (-0.85,-1.95) {$d$};

  \node[figlbl, align=center] at (0.1,-3.35)
    {En $x=0$ los dos denominadores coinciden:
     $\phi\propto\dfrac{q+q'}{\sqrt{d^{2}+y^{2}+z^{2}}}$,};
  \node[figlbl, align=center] at (0.1,-4.25)
    {constante s\'olo si $q'=-q$, y entonces $\phi_0=0$.};

\end{tikzpicture}
